%% ============================================================
%% ReliefPlan AI — Academic Paper
%% Springer Nature two-column style
%% Course: 23CSE473 Neural Networks and Deep Learning, Team 2
%% ============================================================
\documentclass[twocolumn]{article}

%% ── Packages ────────────────────────────────────────────────
\usepackage[T1]{fontenc}
\usepackage[utf8]{inputenc}
\usepackage{lmodern}
\usepackage{microtype}
\usepackage[a4paper,
            top=25mm, bottom=25mm,
            left=18mm, right=18mm,
            columnsep=6mm]{geometry}

\usepackage{amsmath,amssymb}
\usepackage{graphicx}
\usepackage{booktabs}
\usepackage{array}
\usepackage{multirow}
\usepackage{hyperref}
\usepackage{url}
\usepackage{cite}
\usepackage{caption}
\usepackage{subcaption}
\usepackage{xcolor}
\usepackage{algorithm}
\usepackage{algpseudocode}
\usepackage{enumitem}
\usepackage{titlesec}
\usepackage{fancyhdr}
\usepackage{times}        % Springer-style Times font

%% ── Page style ──────────────────────────────────────────────
\pagestyle{fancy}\fancyhf{}
\lhead{\small\textit{ReliefPlan AI — Autonomous Disaster Relief Camp Planning}}
\rhead{\small\thepage}
\renewcommand{\headrulewidth}{0.4pt}

%% ── Section formatting ──────────────────────────────────────
\titleformat{\section}      {\normalfont\large\bfseries}  {\thesection.}  {0.5em}{}
\titleformat{\subsection}   {\normalfont\normalsize\bfseries}{\thesubsection}{0.5em}{}
\titleformat{\subsubsection}{\normalfont\normalsize\itshape}{\thesubsubsection}{0.5em}{}

%% ============================================================
\begin{document}

%% ── Title block (single-column span) ────────────────────────
\twocolumn[{%
\begin{@twocolumnfalse}

\begin{center}
  {\LARGE\bfseries
    ReliefPlan AI: Autonomous Disaster Relief Camp Layout\\[4pt]
    Generation from Satellite Imagery Using DeepLabV3+\\[4pt]
    Terrain Segmentation and Constraint-Based\\[4pt]
    Spatial Optimisation\par}

  \vspace{10pt}

  {\large
    Paarthu Reddy (CB.SC.U4CSE23639)$^{1}$\quad
    Shivam Kharat (CB.SC.U4CSE23445)$^{1}$
  \par}

  \vspace{6pt}

  {\normalsize
    $^{1}$Department of Computer Science and Engineering,
    Amrita Vishwa Vidyapeetham, Coimbatore, India\\[2pt]
    \textit{Faculty Supervisor: S. Jenita Bagyammal}\\
    \textit{Course: 23CSE473 -- Neural Networks and Deep Learning}
  \par}

  \vspace{10pt}
  \hrule
  \vspace{8pt}
\end{center}

%% ── Abstract ────────────────────────────────────────────────
\begin{center}{\large\bfseries Abstract}\end{center}
\begin{quotation}\noindent
Rapid deployment of humanitarian relief camps following natural disasters or armed conflicts
remains a critical, time-sensitive challenge.
Existing manual planning approaches --- dependent on expert field teams and printed handbooks ---
frequently lead to overcrowding, inequitable access to water and sanitation, and suboptimal land
utilisation.
This paper presents \textbf{ReliefPlan AI}, an end-to-end AI-driven pipeline that accepts
satellite or drone imagery of a disaster-affected site and automatically outputs a fully annotated,
standards-compliant camp blueprint.
The pipeline comprises five stages:
(1)~image acquisition,
(2)~terrain semantic segmentation with a custom DeepLabV3+ architecture augmented with Atrous
Spatial Pyramid Pooling (ASPP) at dilation rates $r\in\{6,12,18\}$,
(3)~hydrological flood-buffer analysis,
(4)~a constraint-based spatial optimisation engine that places shelter blocks, hierarchical road
networks, potable water tap stands, sanitation blocks, medical centres, food courts, and water
purification reservoirs while satisfying UNHCR Emergency Settlement Standards, and
(5)~a geo-referenced blueprint renderer.
Experimental results show that the proposed DeepLabV3+ model achieves a mean
Intersection-over-Union (mIoU) of \textbf{44.9\%} and Dice coefficient of \textbf{61.1\%} over
five training epochs, outperforming the U-Net baseline (mIoU 30.3\%, Dice 46.1\%) under identical
conditions, while requiring $8{\times}$ fewer parameters (3.9M vs.\ 31.0M) and $2.4{\times}$ lower
inference latency (394.6\,ms vs.\ 936.6\,ms).
The optimisation engine achieves a \textbf{99.4\%} UNHCR compliance score, maximises the feasible
population on buildable land, and guarantees potable water and sanitation access within every
discrete residential block.
ReliefPlan AI is deployable as a web application, enabling field officers to generate a complete
camp blueprint from raw satellite input in under one minute.
\end{quotation}
\vspace{4pt}

\noindent\textbf{Keywords:}\;
Disaster relief camp planning $\cdot$
Semantic segmentation $\cdot$
DeepLabV3+ $\cdot$
Atrous Spatial Pyramid Pooling $\cdot$
UNHCR standards $\cdot$
Spatial optimisation $\cdot$
Satellite image analysis $\cdot$
Humanitarian AI

\vspace{8pt}
\hrule
\vspace{10pt}
\end{@twocolumnfalse}
}]

%% ============================================================
\section{Introduction}
\label{sec:intro}
%% ============================================================

Natural disasters --- earthquakes, floods, tropical cyclones, and armed conflicts --- displace
millions of people each year~\cite{unhcr_handbook}.
Rapid establishment of safe, sanitary, and accessible temporary shelter is a first-order
humanitarian priority; delays in camp deployment directly translate into excess morbidity and
mortality among the affected population.

Current practice assigns camp design to a small team of experienced humanitarian planners who work
under severe time pressure, using a combination of GIS desktop tools and printed field handbooks
published by organisations such as UNHCR and the Sphere Project~\cite{sphere_handbook}.
The resulting layouts often exhibit avoidable inefficiencies: irregular road networks that impede
ambulance access, shelter blocks placed on flood-prone ground, and sanitation facilities sited
beyond walkable distance from sleeping areas.

\textbf{ReliefPlan AI} addresses these limitations through an automated planning pipeline.
Given only a satellite or drone image and an estimated displaced population, the system:
\begin{enumerate}[leftmargin=*,nosep,topsep=2pt]
  \item Performs pixel-level terrain classification to distinguish buildable land, water bodies,
        dense vegetation, roads, and existing structures.
  \item Enforces a hydrological flood-safety buffer around all detected water bodies.
  \item Runs a constraint-based spatial optimisation algorithm placing shelter pods, arterial and
        feeder roads, water tap stands, sanitation blocks, medical triage centres, food courts,
        administrative gates, and water-purification reservoir tanks --- all satisfying UNHCR
        Emergency Settlement Standards.
  \item Renders a colour-coded camp blueprint overlaid on the original satellite image.
  \item Computes humanitarian metrics: estimated population capacity, land-utilisation rate,
        average walking distance to services, and UNHCR compliance score.
\end{enumerate}

The remainder of the paper is structured as follows.
Section~\ref{sec:related} reviews related academic work.
Section~\ref{sec:methodology} details the methodology.
Section~\ref{sec:experiments} reports experimental results.
Section~\ref{sec:discussion} discusses findings and limitations.
Section~\ref{sec:conclusion} concludes with future directions.

%% ============================================================
\section{Related Work}
\label{sec:related}
%% ============================================================

\subsection{Semantic Segmentation of Overhead Imagery}

Semantic segmentation of aerial imagery has been extensively studied.
Long et al.~\cite{fcn} introduced Fully Convolutional Networks (FCNs), establishing the
encoder--decoder paradigm.
Ronneberger et al.~\cite{unet} extended this with U-Net's symmetric skip connections, which became
the de-facto baseline for remote-sensing segmentation.
Chen et al.~\cite{deeplabv3plus} proposed DeepLabV3+ with Atrous Spatial Pyramid Pooling (ASPP)
to capture multi-scale context without spatial resolution loss --- particularly beneficial for
aerial imagery, where relevant objects span pixel widths from narrow roads to broad water bodies.

Public benchmarks for aerial land-cover segmentation include LoveDA~\cite{loveda},
SpaceNet~\cite{spacenet}, xView~\cite{xview}, and the Sentinel-2 multispectral
archive~\cite{sentinel2}.

\subsection{Facility Layout Optimisation and Reinforcement Learning}

The Facility Layout Problem (FLP) asks how to arrange functional units within a bounded area to
minimise travel cost.
Recent research has applied Deep Reinforcement Learning (DRL) --- Proximal Policy Optimisation
(PPO), Advantage Actor-Critic (A2C), and Deep Q-Networks --- to FLP, demonstrating competitive
arrangements that generalise across differently shaped layouts without
retraining~\cite{flp_drl_2026,flp_generalise_2025}.
The open-source \texttt{gym-flp}~\cite{gymflp} environment standardises DRL benchmarking for FLP
within the OpenAI Gym framework.

\subsection{Automated Humanitarian Settlement Planning}

The intersection of AI and humanitarian settlement planning remains significantly underexplored.
UNHCR and the Sphere Project publish detailed field manuals encoding decades of practical
experience~\cite{unhcr_handbook,sphere_handbook}, but these standards are consumed manually by
planners rather than encoded as machine-readable constraints.
Satellite imagery has been used for rapid damage assessment~\cite{copernicus_em} and shelter
mapping, but no end-to-end system that automatically generates complete camp layouts from raw
imagery exists in the open literature.
ReliefPlan AI bridges both bodies of work, treating humanitarian camp design as a structured layout
optimisation problem with UNHCR/Sphere standards encoded directly as constraint parameters.

%% ============================================================
\section{Methodology}
\label{sec:methodology}
%% ============================================================

\subsection{System Architecture}

The ReliefPlan AI pipeline comprises five sequential modules:
\textbf{Image Acquisition} $\rightarrow$
\textbf{Terrain Segmentation} $\rightarrow$
\textbf{Terrain Analysis \& Flood Buffer} $\rightarrow$
\textbf{Constraint-Based Layout Optimisation} $\rightarrow$
\textbf{Blueprint Renderer}.
Each module addresses a specific gap left by manual planning practice.

\subsection{Terrain Segmentation Model}

\subsubsection{DeepLabV3+ Architecture with ASPP}

The segmentation module implements a custom DeepLabV3+~\cite{deeplabv3plus}.
A lightweight convolutional backbone produces:
\begin{itemize}[leftmargin=*,nosep,topsep=2pt]
  \item Low-level features:
    $F_\mathrm{low}\in\mathbb{R}^{64\times\frac{H}{2}\times\frac{W}{2}}$
  \item High-level features:
    $F_\mathrm{high}\in\mathbb{R}^{256\times\frac{H}{16}\times\frac{W}{16}}$
\end{itemize}

\paragraph{ASPP Module.}
Five parallel branches aggregate multi-scale context from $F_\mathrm{high}$:
\begin{align}
  \mathrm{ASPP}(x) &= \mathcal{P}\!\Bigl(
    f_{1\times1}(x)\;\|\;
    f_{3\times3}^{r=6}(x)\;\|\;
    f_{3\times3}^{r=12}(x)\;\|\; \nonumber\\
  &\qquad\qquad f_{3\times3}^{r=18}(x)\;\|\;
    \mathrm{GAP}(x)
  \Bigr)
  \label{eq:aspp}
\end{align}
where $f_{k\times k}^{r}$ is a dilated convolution with dilation rate $r$, Batch Normalisation,
and ReLU; $\mathrm{GAP}$ is Global Average Pooling upsampled to match $x$; $\|$ denotes channel
concatenation; and $\mathcal{P}$ is a $1\times1$ projection to 256 channels with 0.5 Dropout.

\paragraph{Decoder.}
Low-level features are projected to 48 channels, concatenated with the upsampled ASPP output
($256+48=304$ channels), processed by two $3\times3$ convolutions, and finally upsampled to the
input resolution:
\begin{equation}
  \hat{Y} = \mathrm{Up}\!\left(
    W_c \ast \mathrm{Dec}\!\left[
      \mathrm{ASPP}(F_\mathrm{high})\;\|\;
      \mathcal{P}_l(F_\mathrm{low})
    \right]
  \right)
  \label{eq:decoder}
\end{equation}

\paragraph{Segmentation classes.}
Seven semantic classes: open/buildable land, water body, dense vegetation, road, existing
structure, barren/rocky terrain, background.

\subsubsection{Training Configuration}

\begin{itemize}[leftmargin=*,nosep,topsep=2pt]
  \item Optimiser: Adam, $\mathrm{lr}{=}1\!\times\!10^{-3}$,
        weight decay $1\!\times\!10^{-4}$
  \item Schedule: Cosine Annealing over 5 epochs
  \item Loss: Weighted cross-entropy (class-frequency inverse weights)
  \item Input: $256\times256$; augmentation: horizontal/vertical flip, colour jitter
\end{itemize}

\subsection{Hydrological Flood-Buffer Analysis}

A binary water-body mask $M_W$ is dilated by an elliptical structuring element of radius 11 pixels
($\approx$15--20\,m real-world setback at typical satellite image resolution) to create the
flood-safety exclusion zone $M_F$:
\begin{equation}
  M_F = M_W \oplus E_{r=11}
  \label{eq:flood}
\end{equation}
The buildable mask is updated: $M_B \leftarrow M_B \setminus M_F$.

A water-purification tank is placed at the optimal position $\mathcal{T}^*$ adjacent to the water
body yet outside the exclusion zone:
\begin{equation}
  \mathcal{T}^* = \!\!\argmin_{\substack{p:\,7 \le d_W(p) \le 25 \\ M_F(p)=0}} \!\! \lVert p - c_W \rVert
  \label{eq:tank}
\end{equation}
where $d_W(p)$ is the Euclidean distance from pixel $p$ to the nearest water-body pixel, and $c_W$
is the water-body centroid.

\subsection{Constraint-Based Layout Optimisation Engine}

\subsubsection{Buildable Parcel Detection}

The cleaned buildable mask is analysed by connected-component labelling.
Parcels smaller than $\max(70,\lfloor HW/500 \rfloor)$ pixels are discarded.
Surviving parcels are ranked by area; the largest parcel receives priority placement of the medical
centre, food court, and administrative gate.

\subsubsection{Hierarchical Road Network}

Primary block size: $b = \max(12, \lfloor\min(H,W)/20\rfloor)$;
road width: $\rho = \max(2, \lfloor b/5 \rfloor)$.
An arterial spine runs through the centroid of each parcel along its longer axis.
Perpendicular feeder roads, spaced every
$\Delta = \max(2b+4,\;22)$ pixels, subdivide each parcel into
discrete rectangular residential \emph{boxes}:
\begin{equation}
  R = \bigl\{(r,c)\,\bigm|\,
    c \equiv 0\!\pmod{\Delta},\;(r,c)\in\mathrm{parcel}\bigr\}
  \label{eq:road}
\end{equation}
This guarantees every shelter block is road-accessible.

\subsubsection{Mandatory Per-Box Service Allocation}

For every residential box $\mathcal{B}_i$ formed by the road grid, the engine enforces two
\emph{hard} constraints before shelter packing:
\begin{itemize}[leftmargin=*,nosep,topsep=2pt]
  \item One \textbf{potable water tap stand} at the road-facing edge of the box.
  \item One \textbf{sanitation/latrine block} at the interior boundary corner of the box.
\end{itemize}

\subsubsection{UNHCR Standards as Optimisation Constraints}

Table~\ref{tab:unhcr} lists the UNHCR Emergency Settlement Standards encoded as
constraint parameters in the optimisation engine.

\begin{table}[h]
  \caption{UNHCR standards encoded as constraint parameters~\cite{unhcr_handbook}.}
  \label{tab:unhcr}
  \centering\small
  \setlength{\tabcolsep}{4pt}
  \begin{tabular}{@{}p{4.3cm}r@{}}
    \toprule
    \textbf{Standard} & \textbf{Threshold} \\
    \midrule
    Total camp area per person            & $\ge 45\,\mathrm{m}^2$ \\
    Covered shelter per person            & $\ge 3.5\,\mathrm{m}^2$ \\
    Persons per family tent               & $\approx 5$ \\
    Latrines per persons                  & 1 per 20 \\
    Showers per persons                   & 1 per 50 \\
    Max walk to latrine                   & $\le 50\,\mathrm{m}$ \\
    Water tap stands per persons          & 1 per 80 \\
    Max walk to water point               & $\le 200\,\mathrm{m}$ \\
    Health centre per persons             & 1 per 20\,000 \\
    Food distribution point per persons   & 1 per 5\,000 \\
    \bottomrule
  \end{tabular}
\end{table}

\subsubsection{Shelter Density Optimisation}

After service allocation, empty cells inside each box are packed with family shelter blocks in two
passes:
\begin{itemize}[leftmargin=*,nosep,topsep=2pt]
  \item \textbf{Primary pass}: $b{\times}b$ blocks (4--5 tents/unit), stride $b+1$, occupancy
        threshold $\ge 78\%$.
  \item \textbf{Infill pass}: $b_i{\times}b_i$ blocks ($b_i=\lfloor b/2\rfloor$, 2 tents/unit),
        stride $b_i+1$, threshold $\ge 75\%$.
\end{itemize}
Edge fragments are filled by single-tent infill units at a 70\% threshold.

\subsubsection{Population and Metrics Calculation}

\noindent\textbf{Estimated population capacity:}
\begin{equation}
  P = 5\cdot\bigl(4\,S_\mathrm{prim} + 2\,S_\mathrm{inf}\bigr)
  \label{eq:pop}
\end{equation}

\noindent\textbf{Land-utilisation rate:}
\begin{equation}
  U = \frac{\bigl|\{(r,c):G(r,c)>0\}\bigr|}
           {\bigl|\{(r,c):M_B(r,c)>0\}\bigr|}
      \times 100\%
  \label{eq:util}
\end{equation}

\noindent\textbf{UNHCR Compliance Score} (weighted adherence indicators):
\begin{equation}
  S_\mathrm{UNHCR} = \sum_{k} w_k \cdot \mathbf{1}[c_k \text{ satisfied}], \quad
  \sum_k w_k = 1
  \label{eq:score}
\end{equation}
Individual indicators $c_k$ cover water proximity, latrine ratio, shelter spacing, road
connectivity, flood setback, health coverage, food court coverage, and administrative access.

\subsection{Humanitarian Logistics Estimates}

\noindent\textbf{Daily water requirement} (WHO: 15\,L/person/day):
\begin{equation}
  W_\mathrm{day} = 15\,P \quad [\mathrm{L/day}]
  \label{eq:water}
\end{equation}

\noindent\textbf{Daily caloric provision} (WFP: 2\,100\,kcal/person/day):
\begin{equation}
  F_\mathrm{day} = 2100\,P \quad [\mathrm{kcal/day}]
  \label{eq:food}
\end{equation}

\noindent\textbf{Solar power requirement} (50\,W/tent):
\begin{equation}
  E = 50\cdot S_\mathrm{total} \quad [\mathrm{W}]
  \label{eq:solar}
\end{equation}

\subsection{Web Application}

The system is deployed as a React/Vite frontend communicating with a FastAPI backend over REST.
Users upload satellite imagery, select a population estimate and segmentation model (U-Net or
DeepLabV3+), and receive an annotated blueprint with a full metrics dashboard in under 60 seconds.
Colour coding:
\textit{shelters} (green),
\textit{roads} (grey),
\textit{water taps} (blue),
\textit{latrines} (orange),
\textit{medical} (red),
\textit{food court} (magenta),
\textit{water tank} (teal),
\textit{admin gate} (yellow).

%% ============================================================
\section{Experiments and Results}
\label{sec:experiments}
%% ============================================================

\subsection{Segmentation Model Comparison}

Table~\ref{tab:model} compares DeepLabV3+ (proposed) against U-Net (baseline) trained under
identical conditions: same dataset split, same number of epochs, same augmentation policy.
Model parameters and inference time were measured on a standard CPU (profiling\_results.json).

\begin{table}[h]
  \caption{DeepLabV3+ vs.\ U-Net --- model comparison.}
  \label{tab:model}
  \centering\small
  \setlength{\tabcolsep}{3pt}
  \begin{tabular}{@{}lrrrr@{}}
    \toprule
    \textbf{Model} & \textbf{Params} & \textbf{Inf.\ (ms)} & \textbf{mIoU} & \textbf{Dice} \\
    \midrule
    U-Net (baseline)      & 31.0M & 936.6 & 30.3\% & 46.1\% \\
    DeepLabV3+ (proposed) &  3.9M & 394.6 & \textbf{44.9\%} & \textbf{61.1\%} \\
    \bottomrule
  \end{tabular}
\end{table}

DeepLabV3+ achieves \textbf{14.6\,pp} higher mIoU and \textbf{15.0\,pp} higher Dice coefficient
while using only 12.7\% of the parameters and 42.1\% of the inference time of U-Net.
The performance gap stems from ASPP's multi-scale receptive field, which captures satellite scene
structure across object sizes (narrow roads to broad water bodies) that U-Net's single-scale skip
connections cannot match without substantially more parameters.

\subsection{Training Dynamics}

Table~\ref{tab:training} reports per-epoch validation metrics for DeepLabV3+.

\begin{table}[h]
  \caption{DeepLabV3+ validation performance per epoch.}
  \label{tab:training}
  \centering\small
  \begin{tabular}{@{}rrrr@{}}
    \toprule
    \textbf{Epoch} & \textbf{Val Loss} & \textbf{Val mIoU} & \textbf{Val Dice} \\
    \midrule
    1 & 1.308 & 28.8\% & 42.9\% \\
    2 & 1.293 & 36.4\% & 51.3\% \\
    3 & 1.053 & 42.0\% & 58.1\% \\
    4 & 1.166 & 43.5\% & 59.5\% \\
    5 & \textbf{1.005} & \textbf{44.9\%} & \textbf{61.1\%} \\
    \bottomrule
  \end{tabular}
\end{table}

Both mIoU and Dice coefficient improve monotonically over five epochs, with the largest single-epoch
gain occurring between epochs 1 and 3, consistent with the model learning coarse terrain boundaries
(land vs.\ water vs.\ road) before refining finer class distinctions.

\subsection{Layout Optimisation Results}

On a representative $512{\times}512$ satellite image the optimisation engine produces:

\begin{table}[h]
  \caption{Layout optimisation key metrics.}
  \label{tab:layout}
  \centering\small
  \begin{tabular}{@{}lr@{}}
    \toprule
    \textbf{Metric} & \textbf{Value} \\
    \midrule
    UNHCR Compliance Score           & 99.4\% \\
    Land Utilisation                 & 85--95\% \\
    Avg.\ walk to nearest service    & 34.2\,m \\
    Per-box water tap (guaranteed)   & Yes (100\%) \\
    Per-box latrine block (guaranteed)& Yes (100\%) \\
    Water tank placed near water body & Yes \\
    \bottomrule
  \end{tabular}
\end{table}

These results demonstrate that the engine satisfies all UNHCR hard constraints while maximising
shelter density and minimising service walking distances.

\subsection{Datasets Used}

Training and evaluation used a curated subset drawn from:
SpaceNet~\cite{spacenet} (building/road extraction),
LoveDA~\cite{loveda} (land-cover segmentation across urban and rural scenes),
xView~\cite{xview} (overhead imagery object detection), and
Sentinel-2~\cite{sentinel2} (multispectral satellite archive).
OpenStreetMap annotations supplemented road and building labels.

%% ============================================================
\section{Discussion}
\label{sec:discussion}
%% ============================================================

\subsection{Strengths}

ReliefPlan AI is, to the authors' knowledge, the first system to close the loop from raw satellite
imagery to a fully annotated, standards-compliant camp blueprint.
Key strengths are:
\begin{enumerate}[leftmargin=*,nosep,topsep=2pt]
  \item ASPP's multi-scale receptive field is well-matched to the scale variance inherent in
        satellite imagery.
  \item Hard per-box water-and-latrine constraints guarantee zero residential areas without
        sanitation access, a first-principles improvement over manual methods.
  \item The hydrological flood-buffer and adjacent water-tank placement translate physical
        geospatial reasoning into encoded constraints.
  \item The web interface reduces expert access barriers, enabling any field officer with a
        satellite image to generate a compliant blueprint.
\end{enumerate}

\subsection{Limitations}

\begin{enumerate}[leftmargin=*,nosep,topsep=2pt]
  \item \textit{Limited training data}: mIoU of 44.9\% is below the 70+\% typical of models
        pre-trained on large-scale benchmarks.
        Fine-tuning from a ResNet-50 backbone on full SpaceNet/LoveDA would substantially improve
        accuracy.
  \item \textit{Heuristic optimiser}: the grid-carving strategy is efficient but may not find the
        globally optimal layout for highly irregular terrain shapes.
  \item \textit{Static planning}: the current system does not support dynamic camp expansion as
        displaced population grows.
\end{enumerate}

\subsection{Future Work}

Planned extensions include:
replacing the constraint-based engine with a PPO-trained DRL agent~\cite{flp_drl_2026};
3D terrain modelling; multi-drone mapping; live displacement-data integration;
population-growth prediction for dynamic camp planning; IoT-based resource monitoring; and a
mobile application for on-ground disaster response teams.

%% ============================================================
\section{Conclusion}
\label{sec:conclusion}
%% ============================================================

This paper presented ReliefPlan AI, an end-to-end AI pipeline that autonomously generates disaster
relief camp blueprints from satellite imagery.
The proposed DeepLabV3+ model with ASPP outperforms the U-Net baseline by 14.6\,pp in mIoU,
requiring $8{\times}$ fewer parameters and $2.4{\times}$ lower inference latency.
The constraint-based optimisation engine achieves a 99.4\% UNHCR compliance score and guarantees
water and sanitation access in every residential block.
The system directly supports UN Sustainable Development Goals 3 (Good Health and Well-being),
6 (Clean Water and Sanitation), 11 (Sustainable Cities and Communities), and 16 (Peace, Justice
and Strong Institutions), offering humanitarian organisations a practical tool to accelerate camp
deployment while improving safety and equity outcomes for displaced populations.

\vspace{6pt}
\noindent\textbf{Acknowledgements.}
The authors thank Prof.\ S.\ Jenita Bagyammal for her supervision throughout the project.
Satellite imagery datasets used in this work are publicly available through the sources cited below.

%% ── References ──────────────────────────────────────────────
\begin{thebibliography}{99}

\bibitem{unhcr_handbook}
UNHCR, \textit{Emergency Handbook: Camp Planning Standards (Planned Settlements)},
United Nations High Commissioner for Refugees, Geneva, 2022.
\url{https://emergency.unhcr.org/entry/40977}

\bibitem{sphere_handbook}
Sphere Project, \textit{Humanitarian Charter and Minimum Standards in Humanitarian Response},
4th ed., Practical Action Publishing, Rugby, 2018.

\bibitem{fcn}
J.~Long, E.~Shelhamer, and T.~Darrell,
``Fully convolutional networks for semantic segmentation,''
in \textit{Proc.\ IEEE CVPR}, pp.~3431--3440, 2015.

\bibitem{unet}
O.~Ronneberger, P.~Fischer, and T.~Brox,
``U-Net: Convolutional networks for biomedical image segmentation,''
in \textit{Proc.\ MICCAI}, LNCS vol.~9351, pp.~234--241, 2015.

\bibitem{deeplabv3plus}
L.-C.~Chen, Y.~Zhu, G.~Papandreou, F.~Schroff, and H.~Adam,
``Encoder-decoder with atrous separable convolution for semantic image segmentation,''
in \textit{Proc.\ ECCV}, LNCS vol.~11211, pp.~833--851, 2018.

\bibitem{loveda}
J.~Wang \textit{et al.},
``LoveDA: A remote sensing land-cover dataset for domain adaptive semantic segmentation,''
in \textit{Proc.\ NeurIPS Datasets and Benchmarks Track}, 2021.

\bibitem{spacenet}
A.~Van Etten, D.~Lindenbaum, and T.~M.~Bacastow,
``SpaceNet: A remote sensing dataset and challenge series,''
\textit{arXiv preprint} arXiv:1807.01232, 2018.

\bibitem{xview}
D.~Lam \textit{et al.},
``xView: Objects in context in overhead imagery,''
\textit{arXiv preprint} arXiv:1802.07856, 2018.

\bibitem{sentinel2}
M.~Drusch \textit{et al.},
``Sentinel-2: ESA's optical high-resolution mission for GMES operational services,''
\textit{Remote Sensing of Environment}, vol.~120, pp.~25--36, 2012.

\bibitem{gymflp}
M.~Schweidler \textit{et al.},
``gym-flp: A reinforcement learning environment for facility layout problems,''
\textit{SoftwareX}, 2022.
\url{https://github.com/florianschweidler/gym-flp}

\bibitem{flp_drl_2026}
X.~Zhang and Y.~Liu,
``Optimization of facility layout using deep reinforcement learning,''
\textit{Journal of Computational Design and Engineering},
vol.~13, no.~2, pp.~412--428, 2026.

\bibitem{flp_generalise_2025}
K.~M\"{u}ller and P.~Hoang,
``Investigating the generalizability of facility layout problems using deep reinforcement
learning,''
\textit{Production Engineering (Springer)}, vol.~19, pp.~145--162, 2025.

\bibitem{copernicus_em}
Copernicus Emergency Management Service,
``Rapid mapping and early warning for disasters,''
European Commission Joint Research Centre, 2023.
\url{https://emergency.copernicus.eu}

\end{thebibliography}

\end{document}
